% Chapter 4, section 4: The Laravel worker supervisor
% Draft by Claude, 2026-10-24. Run on QueenMQ 2.6.0 with PHP client 2.5.0 and Laravel 12, in a copy
% of examples/apps/laravel: both engines (the PHP one, and the Rust one, supervisor 0.8.1 built
% from supervisor/), the 300-job burst behind the figure, status, pause, continue, terminate and
% status --check. queen:supervisor-install (which downloads the signed binary) was not run.

\section{The Laravel Worker Supervisor}

QueenMQ is also a queue driver for Laravel. Your jobs stay Laravel jobs, \texttt{dispatch()} does not
change, and \texttt{php artisan queue:work} is still the worker. Install the package, then point
Laravel at the broker in \texttt{.env}:

\begin{shell}
composer require queen-mq/php-client
php artisan vendor:publish --tag=queen-config
\end{shell}

\begin{shell}
QUEUE_CONNECTION=queen
QUEEN_URL=http://127.0.0.1:6632
\end{shell}

The \textbf{supervisor} is what runs those workers. It is a master process that starts ordinary
\texttt{queue:work} processes, keeps the right number of them on each queue, restarts the ones that
exit, and stops them cleanly at a deploy. It does the job of Horizon's master, with no Redis.

\why{a fixed number of workers is too many at night and too few in a burst. A master that reads the
backlog sizes the pool for you, and a clean stop means a deploy never cuts a job in half.}

\queenfig{supervisor-topology}{The supervisor starts the workers; the workers talk to QueenMQ.}

\subsection{Two engines, one configuration}

\begin{pairs}[0.12]{Engine}{What it is}
PHP & Runs straight from the Composer package, with Laravel loaded in its master. Start it with
\texttt{php artisan queen:supervise}. \\
Rust & A small native master: 7.5~MiB of memory in our run. Install it once with
\texttt{php artisan queen:supervisor-install}, then start \texttt{vendor/bin/queen-supervisor}. \\
\end{pairs}

Both read the same \texttt{supervisor} section of \texttt{config/queen.php}, write the same status and
obey the same commands. Start with one and move to the other without touching a pool.

\subsection{How to use it}

The supervisor keeps its state (its status, its lock, the commands you send it) in a private
directory, owned by the user that runs it. It refuses a directory others could write to:

\begin{shell}
install -d -m 0700 /srv/app/queen-supervisor
echo "QUEEN_SUPERVISOR_STATE_DIRECTORY=/srv/app/queen-supervisor" >> .env
\end{shell}

Then describe a pool in \texttt{config/queen.php}. A pool is a set of workers with one configuration,
like a Horizon supervisor. This one serves the queue \texttt{default} with one to eight workers:

\begin{code}
'supervisor' => [
    'supervisors' => [
        'default' => [
            'connection' => 'queen',
            'queues' => ['default'],
            'balance' => 'auto',              // the pool follows the backlog
            'min_processes' => 1,
            'max_processes' => 8,
            'target_jobs_per_process' => 10,  // one worker per 10 waiting jobs
            'fast_scale_up' => true,          // close half the gap at each step
            'timeout' => 60,                  // shorter than retry_after (90 s)
            'tries' => 3,
        ],
    ],
],
\end{code}

Start it, with either engine, under the process manager you already use (systemd, Supervisor,
Kubernetes), and ask it how it is:

\begin{shell}
php artisan queen:supervise                               # the PHP engine
vendor/bin/queen-supervisor --php php --artisan artisan   # or the Rust engine

php artisan queen:supervisor status
+--------+---------+------+-------+---------------+-------+----------------------+
| engine | state   | live | ready | full capacity | pid   | updated              |
+--------+---------+------+-------+---------------+-------+----------------------+
| php    | running | yes  | yes   | yes           | 45516 | 2026-10-24T15:03:55Z |
+--------+---------+------+-------+---------------+-------+----------------------+
\end{shell}

\subsection{What it does with a burst}

We dispatched 300 jobs of 200~ms each onto this pool. Three seconds later the supervisor had seen
290 waiting jobs and grown the pool from one worker to five. It had eight at nine seconds, and the
queue was empty at fifteen. Ten seconds after that it began to give the extra workers back, one at a
time, until one was left. The Rust engine did the same, within a second.

\queenfig{supervisor-run}{A burst of 300 jobs, as the supervisor followed it.}

\begin{options}{Pool option}{Default}{What it does}
balance & auto & \texttt{auto}: the pool follows the backlog. \texttt{simple}: a fixed number of
workers. \texttt{off}: every worker serves the whole ordered list of queues. \\
min\_processes, max\_processes & 1, 10 & The bounds of the pool. \\
strategy & size & \texttt{size}: one worker per \texttt{target\_jobs\_per\_process} waiting jobs.
\texttt{time}: enough workers to clear the backlog within \texttt{target\_clear\_seconds}. \\
balance\_max\_shift, balance\_cooldown & 1, 3~s & How many workers one step adds or removes, and how
often. \texttt{fast\_scale\_up} closes half the gap at each step instead. \\
scale\_down\_delay & 10~s & How long a lower target must hold before workers are given back. \\
timeout, retry\_after & 60~s, 90~s & A job's time limit, and its lease. The supervisor refuses to
start when the timeout is not the shorter one. \\
\end{options}

\subsection{Control it}

\begin{shell}
php artisan queen:supervisor pause           # workers finish their job, take no new ones
php artisan queen:supervisor continue        # and start again
php artisan queen:supervisor terminate       # drain and exit, at every deploy
php artisan queen:supervisor status --check  # a probe: fails unless every pool serves
\end{shell}

While we paused it, 20 new jobs waited in the queue and \texttt{status --check} failed; after
\texttt{continue} they were done in a few seconds. \texttt{terminate} stopped the master and every
worker in two seconds.

\why{a deploy must not cut a job in half, and what runs after it must be the new code.
The terminate command lets each worker finish its job and hands the jobs it had not started back
to the broker. Your process manager then starts a fresh master, which loads the new code.}

\subsection{Four things to know}

\begin{itemize}
\item \textbf{One master per application and consumer group.} A master sizes its pools from the
whole backlog, so two of them would each run the full target. To run several replicas, turn on
\texttt{QUEEN\_SUPERVISOR\_COORDINATION=true} in every one: they then share one target.
\item \textbf{Prefork saves most of the memory.} With \texttt{QUEEN\_SUPERVISOR\_PREFORK=true},
Laravel boots once and every worker is forked from it: in our benchmark, 64 forked workers fit in
183~MiB, where Horizon's took 1,892.
\item \textbf{It refuses an unsafe pool at the start}: a timeout that is not shorter than the lease,
or pools whose maximums do not fit its process budget. A mistake shows at the deploy, not as a job
that runs twice.
\item \textbf{It runs on your worker hosts, never on the broker}, and on Unix only.
\end{itemize}

On the site: \href{https://queenmq.com/integrations/laravel/}{queenmq.com/integrations/laravel}. More
there: \href{https://queenmq.com/integrations/laravel/supervisors/}{worker supervisors},
\href{https://queenmq.com/integrations/laravel/migrate-from-horizon/}{migrate from Horizon}.
